\section{Research Notes and Selected Sources}

This chapter synthesizes foundational principles across classical decision theory, behavioral economics, judgmental forecasting tournaments, and empirical entrepreneurship scholarship. The mathematical and strategic models articulated in this chapter are designed as pedagogical instruments to discipline human intuition; they do not claim that complex commercial realities can be collapsed into simplistic arithmetic certainty.

\begin{enumerate}[leftmargin=14pt,itemsep=3pt,topsep=2pt,parsep=0pt]
  \item \textbf{Frank H. Knight (1921), \textit{Risk, Uncertainty and Profit}.} Boston: Houghton Mifflin. The foundational economic treatise establishing the formal boundary between quantified probabilistic distribution (risk) and unquantifiable ignorance (true uncertainty).

  \item \textbf{Daniel Kahneman and Dan Lovallo (1993), ``Timid Choices and Bold Forecasts: A Cognitive Perspective on Risk Taking.''} \textit{Management Science}, 39(1):17--31. The seminal empirical paper establishing the cognitive mechanisms of the inside view versus the outside view, base-rate neglect, and executive over-optimism.

  \item \textbf{Welton Chang, Eva Chen, Barbara Mellers, and Philip Tetlock (2023), ``Developing Expert Political Judgment: The Impact of Training and Practice on Judgmental Accuracy in Geopolitical Forecasting Tournaments.''} \textit{Judgment and Decision Making}. Demonstrates that structured calibration training and deliberate probabilistic scoring yield durable improvements in subjective accuracy.

  \item \textbf{Douglas W. Hubbard (2014), \textit{How to Measure Anything: Finding the Value of Intangibles in Business}.} Hoboken: John Wiley \& Sons. Specifically the analysis of the Expected Value of Perfect Information (EVPI) and the Expected Value of Sample Information (EVSI) applied to commercial resource allocation.

  \item \textbf{Rita Gunther McGrath and Ian MacMillan (1995), ``Discovery-Driven Planning.''} \textit{Harvard Business Review}, 73(4):44--54. Introduces the methodology of reverse income statements, assumption audits, and milestone-governed learning in high-uncertainty ventures.

  \item \textbf{Jeff Bezos (2015, 2016), \textit{Annual Letters to Amazon Shareholders}.} Outlines the structural operational taxonomy of Type 1 (one-way door, irreversible) versus Type 2 (two-way door, reversible) corporate decision-making and decision velocity.

  \item \textbf{Rita Gunther McGrath (1999), ``Falling Forward: Real Options Reasoning and Entrepreneurial Failure.''} \textit{Academy of Management Review}, 24(1):13--30. Formalizes the application of financial options theory to entrepreneurial experimentation, bounded downside loss, and staged capital deployment.

  \item \textbf{Saras D. Sarasvathy (2001), ``Causation and Effectuation: Toward a Theoretical Shift from Economic Inevitability to Entrepreneurial Contingency.''} \textit{Academy of Management Review}, 26(2):243--263. The landmark paper contrasting causal forecasting models with effectual logic, affordable loss, and stakeholder co-creation.

  \item \textbf{Robert J. Lempert (2019), ``Robust Decision Making (RDM).''} In \textit{Decision Making under Deep Uncertainty}, edited by V. A. W. J. Marchau et al., 23--51. Cham: Springer Open. Examines how institutions design strategies that remain robust across multiple divergent futures rather than optimizing for a single fragile forecast.

  \item \textbf{Daniel Kahneman and Amos Tversky (1979), ``Prospect Theory: An Analysis of Decision under Risk.''} \textit{Econometrica}, 47(2):263--291. The foundational behavioral paper documenting loss aversion, non-linear probability weighting, and reference-point dependence.

  \item \textbf{Daniel Kahneman, Dan Lovallo, and Olivier Sibony (2011), ``Before You Make That Big Decision\dots''} \textit{Harvard Business Review}, 89(6):50--60. Presents a practical 12-question diagnostic checklist for senior executives to audit strategic proposals for cognitive bias prior to capital deployment.

  \item \textbf{Daniel Kahneman and Gary Klein (2009), ``Conditions for Intuitive Expertise: A Failure to Disagree.''} \textit{American Psychologist}, 64(6):515--526. The definitive synthesis detailing the precise environmental conditions (high validity and rapid feedback) required for human intuition to be trustworthy.

  \item \textbf{Gary Klein (2007), ``Performing a Project Premortem.''} \textit{Harvard Business Review}, 85(9):18--19. Introduces prospective hindsight and the premortem protocol to overcome corporate groupthink and uncover lethal failure mechanisms.

  \item \textbf{Paul Graham (2013), ``Do Things that Don't Scale.''} Essays on early-stage entrepreneurship. Emphasizes direct, unscalable manual engagement with customers as the primary method to generate high-fidelity behavioral signal.

  \item \textbf{Eric Ries (2011), \textit{The Lean Startup: How Today's Entrepreneurs Use Continuous Innovation to Create Radically Successful Businesses}.} New York: Crown Business. Formalizes the Build--Measure--Learn feedback loop and the discipline of validated learning.

  \item \textbf{Rita Gunther McGrath and Ian MacMillan (2000), ``Assessing Technology Projects Using Real Options Reasoning.''} \textit{Research-Technology Management}, 43(4):35--49. Applied frameworks for valuing stage-gate technology development investments under technical and commercial volatility.
\end{enumerate}
